Choosing an LLM's compression level per query promises cheaper edge inference, but it assumes that compressed variants remain useful and that the hardware turns fewer parameters into lower cost. We test both assumptions. We port FLAP, a retraining-free structured pruning method, to SmolLM2-1.7B and obtain a ladder of nested pruned variants of one model, stored as small pruning plans over a single copy of the weights. We build a 19{,}242-question pool from eight benchmarks with reproducible splits and a scoring pipeline that checks every result against the evaluation harness. On a Tesla T4, FLAP raises WikiText-2 perplexity from 18.48 to 25.94 at 20\% and 124.62 at 50\% of prunable parameters removed (Wanda-SP: 549.25 at 50\%; Magnitude-SP fails), reproducing the FLAP paper's method ranking at this scale. Pruning 50\% cuts 1024-token prefill latency by 38\% and energy per generated token by 27\%, but batch-1 decode latency stays at 30.4--31.1\,ms per token for every ratio. On 3{,}849 test questions, macro accuracy falls from 62.5\% (unpruned) to 47.3\% (FLAP 20\%), 40.0\% (FLAP 35\%), 32.6\% (FLAP 50\%); several pruned variants collapse to a constant answer on MMLU and BoolQ and fail to finish GSM8K answers. A per-query oracle over the variants reaches 75.8\% versus 64.1\% for the best fixed variant, but once variants that do not beat a constant answer are excluded the gain shrinks to 1.9 points (66.1\%) for this 20/35/50\% ladder. Calibrating FLAP on train-split questions instead of Wikipedia, with no retraining, raises macro accuracy by 2.4, 5.2, 4.3 points at 10/20/35\% (all McNemar $p\le$7.3e-08) and largely removes the constant-answer collapse at 10--20\%. With a ladder of the unpruned model and in-domain-calibrated FLAP at 10/20/35\%, the informative oracle reaches 71.9\% at 239\,ms per question against 64.1\% at 292\,ms for the unpruned model, an upper bound for a realistic router. Per-query compression-level routing therefore depends on in-domain calibration of the cheap tier and on a device that turns compression into cost savings. Laptop-CPU latencies are pending. 
